% Generated draft: requires mathematical and source review.
\documentclass[11pt]{article}
\usepackage[margin=25mm]{geometry}
\usepackage{fontspec}
\setmainfont{Latin Modern Roman}
\usepackage{amsmath,amssymb,mathtools,unicode-math}
\setmathfont{Latin Modern Math}
\usepackage{enumitem,needspace,xurl,hyperref}
\setlength{\parindent}{0pt}
\setlength{\parskip}{5pt}
\setlength{\emergencystretch}{3em}
\newcommand{\N}{\mathbb{N}}
\newcommand{\Z}{\mathbb{Z}}
\newcommand{\Q}{\mathbb{Q}}
\newcommand{\R}{\mathbb{R}}
\newcommand{\C}{\mathbb{C}}
\newcommand{\F}{\mathbb{F}}
\newcommand{\A}{\mathbb{A}}
\newcommand{\PP}{\mathbb{P}}
\newcommand{\E}{\mathbb{E}}
\newcommand{\eps}{\varepsilon}
\newcommand{\dd}{\,\mathrm{d}}
\begin{document}
\section*{UnsolvedMath ID: 11100024}
\textbf{Source problem code:} catalogue entry
\subsection{Definitions and mathematical statement}
For a simplicial set \(X\), Hovey tentatively proposes a geometric cocycle model for complex cobordism \(MU^*(|X|)\): place manifolds over vertices, bordisms over edges and successive bordisms-between-bordisms over higher simplices, with compatible orientations/face data and an appropriate equivalence relation. He expressly says the higher-bordism meaning and equivalence are unknown to him, and recalls a possibly published description without author or citation.

A precise version must specify grading and tangential structure (complex cobordism requires stably complex data, not merely an unspecified ordinary orientation), coherence on all simplex faces/degeneracies, stabilization, addition/inverses and equivalence. Its required comparison would be natural identifications
\[\mathcal C^q(X)\cong MU^q(|X|)=[\Sigma^\infty|X|_+,\Sigma^q MU]\]
for the proposed geometric classes \(\mathcal C^q(X)\). The notation \(\mathcal C^q\) is a placeholder for the missing construction. A naive fixed-dimensional manifold at every vertex does not itself specify all positive/negative cohomological degrees or homotopy coherence. Manifolds with corners or other higher-bordism models would need explicit definitions rather than an informal iteration.

The item then asks how complex \(K\)-cohomology is obtained from this description. Hovey speculates that bordisms-between-bordisms become identities and reproduce the cocycle condition for vector bundles; he does not assert that this truncation has been proved to yield \(K\)-theory. A complete construction must define the modification and its natural comparison with the selected \(K\)-cohomology theory, including virtual/stable data where needed. Finally it asks how to obtain elliptic cohomology in the same framework; the elliptic theory, geometric input and comparison are not specified, so no unique spectrum or model is silently chosen.

The research summary mentions Pontryagin--Thom/cobordism and iterated or Baas--Sullivan models as published context, and claims a \(K\)-theory interpretation while leaving the elliptic refinement open. It supplies no exact citation or theorem identifying those constructions with the particular simplex-wise model. Those reports do not remove the source's explicit uncertainty or establish the proposed higher-cocycle truncation.
\subsection{Short English statement}
Make the tentative simplex-by-simplex higher-bordism model of \(MU\)-cohomology precise, then explain its \(K\)-theory and elliptic-cohomology versions. The source does not define higher bordisms or their equivalence and treats the \(K\)-theory mechanism as speculation.
\subsection{Sources}
\begin{itemize}
\item Local UnsolvedMath registry, permanent ID \(11100024\), exact source code \texttt{catalogue entry}.
\item Mark Hovey, \emph{Hovey's Algebraic Topology Problem List}, Elliptic cohomology, item 4, surviving mirror: \url{https://www-users.cse.umn.edu/~tlawson/hovey/}. Registry access 2026-07-29; html-ordered-list extraction.
\item The source cannot identify the recalled publication. The summary names Pontryagin--Thom and Baas--Sullivan-type descriptions without a precise bibliographic match or proof of the stated truncation.
\end{itemize}
\end{document}
